\section{第42届复赛第5题：约瑟夫森结的相位动力学}
\subsection{题干}
题目以超导体--薄绝缘层--超导体结构为对象，给出约瑟夫森关系
\[j=j_0\sin\delta,\qquad \frac{\partial\delta}{\partial t}=\frac{2e}{\hbar}U.\]
依次讨论恒定电压下的交流约瑟夫森效应、外磁场造成的相位空间分布和自生磁场，以及 $J$、$R$、$C$ 并联模型中的相位势阱与等离子体频率。
\begin{figure}[htbp]\centering\begin{tikzpicture}[>=Stealth,scale=.9]
  \fill[gray!30] (-.25,-1.35) rectangle (.25,1.35);
  \draw[thick] (-2,-1.35) rectangle (-.25,1.35);
  \draw[thick] (.25,-1.35) rectangle (2,1.35);
  \draw[dashed] (-.7,-1.5)--(-.7,1.5) (.7,-1.5)--(.7,1.5);
  \draw[->,thick] (-2.9,0)--(-2,0) node[midway,above] {$I$};
  \draw[->,thick] (2,0)--(2.9,0) node[midway,above] {$I$};
  \draw[->] (0,-1.8)--(0,1.85) node[above] {$y$};
  \draw[->] (0,0)--(2.7,0) node[below right] {$x$};
  \draw[<->] (-.25,-1.08)--(.25,-1.08) node[midway,below] {$d$};
  \draw[<->] (-.7,-.82)--(-.25,-.82) node[midway,below] {$\lambda$};
  \draw[<->] (.25,-.82)--(.7,-.82) node[midway,below] {$\lambda$};
  \node[rotate=90,align=center] at (-1.12,.25) {超导体\\1};
  \node[font=\small] at (0,-1.62) {薄绝缘层};
  \node[rotate=90,align=center] at (1.12,.25) {超导体\\2};
  \fill (0,0) circle (1.2pt) node[above left] {$O$};
  \node[below] at (0,-2.25) {原题图5a：约瑟夫森结横截面};
\end{tikzpicture}
\caption{第42届复赛第5题原题图5a的TikZ重绘。}\end{figure}
\begin{figure}[htbp]\centering\begin{tikzpicture}[>=Stealth,scale=.9]
  \coordinate (L) at (-2,0); \coordinate (R) at (2,0);
  \draw[->,thick] (-3,0)--(L) node[midway,above] {$I$};
  \draw[->,thick] (R)--(3,0);
  \draw[thick] (L)--(R);
  \draw[thick] (L)--(-2,1.35)--(-.55,1.35) rectangle (.55,1.75)--(2,1.35)--(R);
  \draw[thick] (L)--(-2,-1.35)--(-.35,-1.35);
  \draw[thick] (-.35,-1.75)--(-.35,-.95) (.35,-1.75)--(.35,-.95);
  \draw[thick] (.35,-1.35)--(2,-1.35)--(R);
  \node[above] at (0,1.75) {$R$};
  \node[below] at (0,-1.75) {$C$};
  \node[above] at (0,.05) {$J$};
  \draw[thick] (-.18,-.2)--(.18,.2) (-.18,.2)--(.18,-.2);
  \fill (L) circle (1.3pt) (R) circle (1.3pt);
  \node[below] at (0,-2.35) {原题图5b：$J$、$R$、$C$ 并联等效电路};
\end{tikzpicture}
\caption{第42届复赛第5题原题图5b的TikZ重绘。}\end{figure}
\subsection{参考解答整理}
恒定电压 $U_0$ 下，
\[\delta(t)=\frac{2eU_0}{\hbar}t+\delta(0),\qquad j(t)=j_0\sin\!\left(\frac{2eU_0}{\hbar}t+\delta(0)\right),\qquad T=\frac{\pi\hbar}{eU_0}.\]
若外场 $B_z=B_0$ 且 $U=0$，相位沿结宽的变化为
\[\delta(y)=\delta(0)+\frac{2e(d+2\lambda)B_0}{\hbar}y.\]
积分电流密度可得总电流的干涉因子；非零的零点外场满足
\[B_0=\frac{n\pi\hbar}{e(d+2\lambda)Y},\qquad n=\pm1,\pm2,\ldots\]
计入自生磁场时，
\[\frac{\mathrm d^2\delta}{\mathrm dy^2}=\frac{1}{\lambda_J^2}\sin\delta,\qquad
\lambda_J=\sqrt{\frac{\hbar}{2e\mu_0j_0(d+2\lambda)}}.\]
在RCSJ模型中，三支路电流相加给出
\[C\frac{\hbar}{2e}\frac{\mathrm d^2\delta}{\mathrm dt^2}+\frac{\hbar}{2eR}\frac{\mathrm d\delta}{\mathrm dt}+I_0\sin\delta=I.\]
忽略电阻支路时，等效势可取
\[V(\delta)=-\frac{2e}{\hbar C}\left(I\delta+I_0\cos\delta\right)+\text{常量}.\]
势阱存在条件为 $|I|<I_0$，稳定底部为 $\delta_0=\arcsin(I/I_0)+2n\pi$；小振动满足
\[\omega^2=\frac{2eI_0}{\hbar C}\sqrt{1-\left(\frac{I}{I_0}\right)^2}.\]
最大等离子体频率对应的传播速率比值为
\[\frac{\bar c}{c}=\sqrt{\frac{d}{\epsilon_r(d+2\lambda)}}=0.055.\]
\subsection{转写依据}
题干依据 \texttt{CPHO42-Q05\_transcription\_v1.md}；公式与分值依据答案第14至16页和 \texttt{CPHO42-Q05\_formula\_check\_v1.csv}。图5a、图5b均由原题图重绘，不嵌入原始页面截图。
